%--------------------------------------------------------------------
% -------------------------------------------------------------------
% TRABALHO DE CONCLUSÃO DE CURSO.
% OTIMIZAÇÕES: Prof. Jonatha Costa
% AUTOR: JOHN DAYTON CHAVES CUSTÓDIO
% -------------------------------------------------------------------
% -------------------------------------------------------------------
% INSTITUTO FEDERAL DE EDUCAÇÃO, CIÊNCIA E TNOLOGIA DO CEARÁ - IFCE
% Departamento de Indústria
% Campus Fortaleza
%%%%%%%%%%%%%%%%%%%%%%%%%%%%%%%%%%%%%%%%%%%%%%%%%%%%%%%%%%%%%%%%%%%%%
%%%%%%%%%%%%%%%%%%%%%% Modelo de Monografia %%%%%%%%%%%%%%%%%%%%%%%%%
%%%%%%%%%%%%%%%%%%%%%%%%%%%%%%%%%%%%%%%%%%%%%%%%%%%%%%%%%%%%%%%%%%%%%
%====================================================================
%================== Definição do documento ==========================
%====================================================================
% \documentclass[conference,harvard,brazil,english]{sbatex}
 \documentclass[
 	a4paper,	% Formato do papel utilizado no trabalho.
 	12pt,		% Tamanho da letra utilizada.
 	oneside,	% Para utilização apenas de um lado da folha. Caso contrario twoside. oneside
    chapter=TITLE,	% Títulos em caixa alta.
 	brazil,		% Idioma padrão do trabalho.
	sumario=abnt-6027-2012  % Formata o sumário conforme as recomendações da ABNT NBR 6027:2012
 	]{abntex2}	% Indicação da utilização do capote abntex2.

%====================================================================
%======================== Pacotes Básicos ===========================
%====================================================================
\usepackage{lmodern}		% Utilização da fonte Latin Modern.
\usepackage[T1]{fontenc}	% Seleção de códigos de fonte.
\usepackage[utf8]{inputenc}	% Para a utilização dos acentos na língua Portuguesa.
\usepackage{indentfirst}	% Identação do primeiro paragrafo de cada seção
%\usepackage{nomencl}		
\usepackage{color}			% Pacote de utilização das cores.
\usepackage{graphicx}		% Pacote para utilização dos gráficos.
\usepackage{microtype}		% Pacote de melhorias de justificação
\usepackage{mathptmx}		% Pacote para incluir a fonte Times New Roman
\usepackage[small]{caption} % Instrução para reduzir a fonte do pacote caption
\usepackage{setspace}		% Permitir o ajusto entre linhas

%====================================================================
%========= Pacotes adicionais, usados apenas para abntx2 ============
%====================================================================
%--------------------------------------------------------------------
% Caso o software LaTeX não separre de forma correta a palavra, pode ser feito manualmente, como a seguir.
\hyphenation{com-pu-ta-dor}		% Separação correta da palavra.
\hyphenation{in-fan-tô-me-tro}
% -------------------------------------------------------------------


%====================================================================
%======================= Pacotes de citações ========================
%====================================================================
%\usepackage[brazilian,hyperpageref]{backref}	% Para páginas com as citações no arquivo bibl
\usepackage[alf,
abnt-full-initials=no,
abnt-emphasize=bf,								% Modificação para texto em negrito.
abnt-and-type=e,
abnt-url-package=url,
abnt-repeated-author-omit=yes]{abntex2cite}		% Citações padrão ABNT



%====================================================================
%========== Espaçamentos entre linhas, parágrafos e margens =========
%====================================================================
% Alterando as margens padrões:
\setlrmarginsandblock{3cm}{2cm}{*}
\setulmarginsandblock{3cm}{2cm}{*}
\checkandfixthelayout

% Define o tamanho do parágrafo:
\setlength{\parindent}{0mm}		% Define parágrafo de 0cm

% Controle do espaçamento entre um parágrafo e outro:
\setlength{\parskip}{0mm}

% Espaçamento Simples
%\SingleSpace


\begin{document}

% Numeração em romanos para páginas iniciais (capa, agradecimentos, etc)
%\pagestyle{empty}
%\pagenumbering{roman}

% -------------------------------------------------------------------
% Determinando as Fontes do documento para os títulos, subtítulos e partes.
% -------------------------------------------------------------------
\renewcommand{\ABNTEXchapterfont}{\fontseries{b}\selectfont}	% Modificar letra do capítulo para negrito
\renewcommand{\ABNTEXchapterfontsize}{\normalsize} 	% Modifica o tamanho da fonte para os capítulos

\renewcommand{\ABNTEXsectionfontsize}{\normalsize} 	% Modifica o tamanho da fonte da seção para normal (12tp)

\renewcommand{\ABNTEXsubsectionfont}{\fontseries{b}\fontshape{it}\selectfont}	% Modificar letra da subseção para negrito em itálico
\renewcommand{\ABNTEXsubsectionfontsize}{\normalsize} 	% Modifica o tamanho da fonte da subseção para normal

\renewcommand{\ABNTEXsubsubsectionfont}{\fontshape{it}\selectfont}	% Modificar letra da subsubseção para itálico

\renewcommand{\ABNTEXsubsubsubsectionfont}{\fontfamily{u}\fontshape{it}\selectfont}	% Modificar letra da subsubsubseção para itálico
%--------------------------------------------------------------------

% Páginas iniciais

%%%%%%%%%%%%%%%%%%%%%%%%%%%%%%%%%%%%%%%%%%%%%%%%%%%%%%%%%%%%%%%%%%%%%
%%%%%%%%%%%%%%%%%%%%%%%%%% CAPA %%%%%%%%%%%%%%%%%%%%%%%%%%%%%%%%%%%%%
%%%%%%%%%%%%%%%%%%%%%%%%%%%%%%%%%%%%%%%%%%%%%%%%%%%%%%%%%%%%%%%%%%%%%

\vspace{3cm}

% Inicio de espaçamento de 1,5
\OnehalfSpace
\begin{center}
\includegraphics[scale=0.7]{Figuras/logo_ifce.png}
\end{center}
\vspace{0.5cm}
\begin{center}

\fontsize{12}{12}
\noindent
\textbf{INSTITUTO FEDERAL DE EDUCAÇÃO, CIÊNCIA E TECNOLOGIA DO CEARÁ - IFCE }\\
\textbf{\textit{CAMPUS} MARACANAÚ}\\
\vspace{0.75cm}
\textbf{CURSO DE BACHARELADO EM ENGENHARIA DE CONTROLE E AUTOMAÇÃO}

\vspace{2cm}
\textbf{INSERIR O SEU NOME}\\
\end{center}

\vspace{4cm}
\begin{center}
\textbf{INSERIR O TÍTULO DO TRABALHO}\\

\vspace{7cm}
\textbf{FORTALEZA}\\
\textbf{\the\year{}}
\end{center}
%%%%%%%%%%%%%%%%%%%%%%%%%%%%%%%%%%%%%%%%%%%%%%%%%%%%%%%%%%%%%%%%%%%%%
%%%%%%%%%%%%%%%%%%%%%% FOLHA DE ROSTO %%%%%%%%%%%%%%%%%%%%%%%%%%%%%%%
%%%%%%%%%%%%%%%%%%%%%%%%%%%%%%%%%%%%%%%%%%%%%%%%%%%%%%%%%%%%%%%%%%%%%

\fontsize{12}{12}							% Tamanho da fonte desejada
\vspace{3cm}								% Espaço vertical de 3cm
\begin{center}								% Inicio da centralização do texto
\textbf{INSERIR O SEU NOME}\\				% Iserção do nome do aluno
\vspace{8cm}								% Espaço vertical de 8cm
							
\textbf{INSERIR O TITULO DO TRABALHO}\\		% Inserção do título do trabalho
\vspace{3cm}								% Espaço vertical de 3 cm
\end{center}								% Final de centralização

\begin{flushright}
\begin{minipage}[15cm] {11cm}
Trabalho de conclusão de curso, na forma de (Tipo do Trabalho), apresentado ao curso de (Bacharelado em Engenharia de Mecatrônica ou Mecatrônica Industrial) do Instituto Federal de Educação, Ciência e Tecnologia do Ceará -- IFCE -- Campus Fortaleza, como requisito para obtenção do Título de Graduação.

\vspace{1cm}
\noindent
Orientador: Prof. Dr. Nome 1\\
Co-orientador: Prof. Msc. Nome 2
\end{minipage}
\end{flushright}
\vspace{6cm}
\begin{center}
FORTALEZA\\
\the\year{}
\end{center}

% Inicio de numeração das páginas. Centralizado e embaixo.
%\pagestyle{plain}
\pagenumbering{arabic}	% Inicia a numeração normal
\pagestyle{empty}		% Não permite a visualização da numeração

%%%%%%%%%%%%%%%%%%%%%%%%%%%%%%%%%%%%%%%%%%%%%%%%%%%%%%%%%%%%%%%%%%%%%
%%%%%%%%%%%%%%%%%%%%%%% FOLHA DE APROVAÇÃO %%%%%%%%%%%%%%%%%%%%%%%%%%
%%%%%%%%%%%%%%%%%%%%%%%%%%%%%%%%%%%%%%%%%%%%%%%%%%%%%%%%%%%%%%%%%%%%%

\begin{center}
\vspace{3cm}
\fontsize{12}{12}
\textbf{INSERIR O SEU NOME}\\
\vspace{1cm}
\textbf{INSERIR O TÍTULO DO TRABALHO}
\vspace{3cm}
\begin{flushright}
\begin{minipage}[15cm] {11cm}
Trabalho de conclusão de curso, na forma de Monografia, apresentado ao curso de Bacharelado em Engenharia de Controle e Automação do Instituto Federal de Educação, Ciência e Tecnologia do Ceará -- IFCE -- \textit{Campus} Maracanaú, como requisito para obtenção do Engenheiro de Controle e Automação, aprovado pela comissão formada pelos professores:

\end{minipage}
\end{flushright}

\vspace{1.5cm}
\begin{flushleft}
\textbf{Aprovado em}\hspace{4pt} \_\_\_\_/\_\_\_\_/\_\_\_\_\_\_\_
\end{flushleft}
\vspace{1.5cm}
\textbf{BANCA EXAMINADORA}

\assinatura{Prof. Dr. Jonatha Rodrigues da Costa \\ Departamento de Indústria, \textit{Campus} Maracanaú (IFCE)}

\vspace{0.8cm}

\assinatura{Prof. Me. Beltrano de Tal \\ Departamento de Indústria, Campus Fortaleza (IFCE)}

\vspace{0.8cm}
\assinatura{Prof. Outro de Tal\\ Departamento de Indústria, Campus Fortaleza (IFCE)}

\end{center}

% Inicio de espaçamento de 1,5
\OnehalfSpace

%%%%%%%%%%%%%%%%%%%%%%%%%%%%%%%%%%%%%%%%%%%%%%%%%%%%%%%%%%%%%%%%%%%%%
%%%%%%%%%%%%%%%%%%%%%%%% DEDICATÓRIA %%%%%%%%%%%%%%%%%%%%%%%%%%%%%%%%
%%%%%%%%%%%%%%%%%%%%%%%%%%%%%%%%%%%%%%%%%%%%%%%%%%%%%%%%%%%%%%%%%%%%%

%====================================================================
% DEDICATÓRIA
\begin{dedicatoria}

	\vspace*{\fill}
    \begin{flushright}
    \begin{minipage}[15cm] {11cm}
    Inserir sua dedicatória aqui...
  \end{minipage}
  \end{flushright}
  \vspace{3cm}

\end{dedicatoria}
%====================================================================
%%%%%%%%%%%%%%%%%%%%%%%%%%%%%%%%%%%%%%%%%%%%%%%%%%%%%%%%%%%%%%%%%%%%%
%%%%%%%%%%%%%%%%%% Folha de Agradecimentos %%%%%%%%%%%%%%%%%%%%%%%%%%
%%%%%%%%%%%%%%%%%%%%%%%%%%%%%%%%%%%%%%%%%%%%%%%%%%%%%%%%%%%%%%%%%%%%%

% ------------------------------------------------------------------- 
% Espaçamentos entre linhas e parágrafos 
% -------------------------------------------------------------------
% O tamanho do parágrafo é dado por:
\setlength{\parindent}{1.25cm}
% Controle do espaçamento entre um parágrafo e outro:
\setlength{\parskip}{0.2cm}
% Inicio de espaçamento de 1,5
\OnehalfSpace
%--------------------------------------------------------------------

%====================================================================
% AGRADECIMENTOS
\begin{agradecimentos}
Meus Agradecimentos aos...

\end{agradecimentos}
%====================================================================
%%%%%%%%%%%%%%%%%%%%%%%%%%%%%%%%%%%%%%%%%%%%%%%%%%%%%%%%%%%%%%%%%%%%%
%%%%%%%%%%%%%%%%%%%%%%%%% EPÍGRAFE %%%%%%%%%%%%%%%%%%%%%%%%%%%%%%%%%%
%%%%%%%%%%%%%%%%%%%%%%%%%%%%%%%%%%%%%%%%%%%%%%%%%%%%%%%%%%%%%%%%%%%%%

\begin{epigrafe}	% Inicio do comando epígrafe
  \vspace*{\fill}	% Espaço para deslocar o testo ao final da página
  \begin{flushright}
  \begin{minipage}[15cm] {11cm}
  ``Texto em que o autor apresenta uma citação, seguida de indicação de autoria, relacionada com a matéria tratada no corpo do trabalho.''
	\flushright{(Nome do autor do texto.)}\\
  \end{minipage}
  \end{flushright}
  
\end{epigrafe}

% ------------------------------------------------------------------- 
% Espaçamentos entre linhas e parágrafos 
% -------------------------------------------------------------------
% O tamanho do parágrafo é dado por:
\setlength{\parindent}{1.25cm}
% Controle do espaçamento entre um parágrafo e outro:
\setlength{\parskip}{0.2cm}  % tente também \onelineskip
%--------------------------------------------------------------------

%%%%%%%%%%%%%%%%%%%%%%%%%%%%%%%%%%%%%%%%%%%%%%%%%%%%%%%%%%%%%%%%%%%%%
%%%%%%%%%%%%%%%%%%%%%%%%% Resumo %%%%%%%%%%%%%%%%%%%%%%%%%%%%%%%%%%%%
%%%%%%%%%%%%%%%%%%%%%%%%%%%%%%%%%%%%%%%%%%%%%%%%%%%%%%%%%%%%%%%%%%%%%

%====================================================================
% RESUMO
\setlength{\absparsep}{18pt} % ajusta o espaçamento dos parágrafos do resumo
\begin{resumo}

Indicações dos pontos mais relevantes em um texto, passando uma rápida e clara visualização dos conteúdos contidos no trabalho. Deverá possuir palavras-chaves logo abaixo da última linha de texto. De preferência utiliza-se aproximadamente três palavras-chaves para representar o conteúdo do trabalho. Devendo estar em português, separadas por ponto, com a primeira letra em maiúsculo e o restante em minúsculo.

 \textbf{Palavras-chave}: modelo de artigo. resumo. template.   
\end{resumo}

%==================================================================
%%%%%%%%%%%%%%%%%%%%%%%%%%%%%%%%%%%%%%%%%%%%%%%%%%%%%%%%%%%%%%%%%%%%%
%%%%%%%%%%%%%%%%%%%%%%%% ABSTRACT %%%%%%%%%%%%%%%%%%%%%%%%%%%%%%%%%%%
%%%%%%%%%%%%%%%%%%%%%%%%%%%%%%%%%%%%%%%%%%%%%%%%%%%%%%%%%%%%%%%%%%%%%


%==================================================================
% ABSTRACT
\setlength{\absparsep}{18pt} % ajusta o espaçamento dos parágrafos do abstract
\begin{resumo}[Abstract]
 \begin{otherlanguage*}{english}

   Concise presentation of the relevant points of a text, providing a quick and clear view of the content and conclusions of the work.


   \textbf{Keywords}: model of article. abstract. template.
 \end{otherlanguage*}
\end{resumo}

% ===================================================================

%%%%%%%%%%%%%%%%%%%%%%%%%%%%%%%%%%%%%%%%%%%%%%%%%%%%%%%%%%%%%%%%%%%%%
%% Lista de figuras, lista de Tabelas, Lista de Símbolos e Sumário %%
%%%%%%%%%%%%%%%%%%%%%%%%%%%%%%%%%%%%%%%%%%%%%%%%%%%%%%%%%%%%%%%%%%%%%


%====================================================================
% LISTA DE ILUSTRAÇÃO
\pdfbookmark[0]{\listfigurename}{lof}

%\renewcommand\listfigurename{\normalsize \textbf{LISTA DE FIGURAS}}
%\renewcommand*{\cftfigurename}{Figure\space} % Essa é a linha de código que irá modificar o nome de Figura para Figure. ;)

\listoffigures*
\cleardoublepage
%====================================================================

%====================================================================
% LISTA DE TABELAS
\pdfbookmark[0]{\listtablename}{lot}
\listoftables*
\cleardoublepage
%====================================================================


%====================================================================
% LISTA DE SIGLAS
\begin{siglas}
\item[IFCE] Instituto Federal de Ciência, Tecnologia e Educação do Ceará
\item[UECE] Universidade Estadual do Ceará
\item[UFC] Universidade Federal do Ceará
\end{siglas}
\cleardoublepage
%====================================================================


%====================================================================
% lISTA DE SÍMBOLOS
\begin{simbolos}
  \item[$ \alpha $] Letra grega Alfa
  \item[$ \beta $] Letra grega Beta
  \item[$ \Gamma $] Letra grega Gama
  \item[$ \Lambda $] Lambda
  \item[$ \zeta $] Letra grega minúscula zeta
  \item[$ \in $] Pertence
\end{simbolos}
%====================================================================

%====================================================================
% SUMÁRIO
%\pdfbookmark[0]{\contentsname}{toc}
\renewcommand{\tocheadstart}{\ABNTEXchapterfont}
\tableofcontents*
%\cleardoublepage
%====================================================================

\SingleSpace
\textual
%\pagestyle{myheadings}		% Permite a visualização da numeração das páginas.
%\pagestyle{headings}		% Permite a visualização da numeração das páginas.
%\pagestyle{plain}

\textcolor{blue}{
\chapter{Prefácio}
Este \textit{template} foi reescrito pelo Prof. Dr. Jonatha Rodrigues da Costa (JRC) a partir do modelo original do IFCE \textit{campus} Fortaleza. Há orientações específicas para cada elemento do texto, fundamentadas em sua experiência como docente, pesquisador e orientador de Trabalhos de Conclusão de Curso na área de Engenharia.

Dessa forma, ao longo deste documento serão apresentadas recomendações metodológicas, técnicas e científicas que refletem as expectativas do orientador quanto à organização do trabalho, ao rigor metodológico, à profundidade técnica e à qualidade da escrita científica.
%

\section{Estruturação}
Orienta-se que o TCC seja escrito de forma clara, técnica e estruturada como segue:

\begin{enumerate}
	\item Introdução
	\begin{enumerate}
		\item Contextualização
		\item Objetivos
		\item Contribuições
		\item Publicações
		\item Organização do trabalho
	\end{enumerate}
	\item Estado da arte
	\begin{enumerate}
		\item Considerações iniciais
		\item Fronteira do conhecimento nos assuntos envolvidos
		\item Visão holística do estado da arte 
		\item Considerações finais
	\end{enumerate}
	\item Fundamentação teórica
	\begin{enumerate}
		\item Considerações iniciais
		\item Conteúdo seminal correlato ao assunto do TCC 
		\item Considerações finais
	\end{enumerate}	
	\item Solução proposta 
	\begin{enumerate}
		\item Arquitetura holística do sistema
		\item Modelo matemático do sistema
		\item Fluxograma do sistema
		\item Considerações finais
	\end{enumerate}
	\item Simulações, Análises e Discussão de resultados
	\begin{enumerate}
		\item Considerações iniciais
		\item Cenários, métodos e parâmetros de simulação/teste/coleta de dados
		\item Apresentação de discussão de resultados: tabelas, gráficos, quadros
		\item Considerações finais
	\end{enumerate}
	\item Conclusões e trabalhos futuros
	\item Referências
	\item Apêndices
	\item Anexos
\end{enumerate}


Esta estruturação garante um mínimo de qualidade técnico-científica a ser aplicada durante a produção de TCC do concluinte de engenharia.
Isto visa, sobre tudo, maximizar o tempo do discente e orientador, reduzindo riscos de ir à defesa com material insuficiente para tanto. Note que este documento está estruturado em \LaTeX\ contendo arquivos \texttt{.tex} já organizados conforme estrutura supracitada. Além disso, há pastas específicas destinadas a organizar o conteúdo em elementos \textbf{pré-textuais}, \textbf{textuais} e \textbf{pós-textuais}.
Desse forma, o concluinte redator do texto deve alterar os arquivos, ou apenas comentar aqueles não aplicáveis ao seu TCC.


% =========================================================
\section{Estrutura lógica e princípio orientador do TCC}
% =========================================================

Um TCC científico deve ser construído como uma sequência lógica e integrada,
na qual cada capítulo responde a uma questão fundamental da pesquisa:

\begin{center}
	\begin{tabular}{>{\bfseries}p{4cm}p{9cm}}
		\toprule
		\textbf{Etapa} & \textbf{Questão central} \\
		\midrule
		Problema & Por que o problema importa? \\
		Estado da Arte & O que já foi feito e o que ainda falta? \\
		Lacuna & O que permanece sem solução? \\
		Objetivo & O que este trabalho pretende realizar? \\
		Fundamentação & Que conhecimento sustenta a pesquisa? \\
		Solução & O que foi desenvolvido? \\
		Metodologia & Como a solução foi avaliada? \\
		Resultados & O que foi observado? \\
		Discussão & O que os resultados significam? \\
		Conclusão & O que foi efetivamente demonstrado? \\
		Trabalhos futuros & O que ainda precisa ser investigado? \\
		\bottomrule
	\end{tabular}
\end{center}

\textbf{
Note que este fluxo deve guiar também a produção da  apresentação (\textit{slildes}) para o defesa.
}

\textbf{Princípio orientador:} no corpo do texto, cada seção deve responder perguntas, cada afirmação relevante deve possuir
evidência, cada resultado deve ser interpretado e cada conclusão deve estar
sustentada pelos resultados apresentados. Essa estrutura garante que o TCC seja mais do que uma descrição de atividades:
ele deve constituir uma \textbf{argumentação científica coerente, fundamentada
	em evidências e capaz de demonstrar uma contribuição}.

Note que uma banca de avaliação pode sempre ser mais rigorosa do que aquilo que se imagina. Portanto, ações proativas mitigam frustrações.

A seguir serão apresentas perguntas-chave para cada capítulo do TCC.


}
\newpage

% =========================================================
\chapter{Introdução}
% =========================================================

\textcolor{blue}{A introdução deve conduzir o leitor do contexto geral do problema até a
delimitação precisa da pesquisa}. Ao final do capítulo, o leitor deve
compreender o problema investigado, sua relevância, os objetivos do trabalho,
suas contribuições e a organização do documento.

% ---------------------------------------------------------
\section{Contextualização}
% ---------------------------------------------------------

O objetivo desta seção é conduzir o leitor do contexto amplo para o problema
científico ou tecnológico específico.

\begin{enumerate}
	\item Qual é o contexto científico, tecnológico, industrial ou social no qual o trabalho está inserido?
	\item Qual fenômeno, sistema, tecnologia ou problema constitui o objeto central da pesquisa?
	\item Por que esse tema é relevante atualmente?
	\item Qual problema concreto ou científico ainda precisa ser resolvido?
	\item Quais são as consequências técnicas, científicas, econômicas ou sociais desse problema?
	\item O que já foi desenvolvido para enfrentar esse problema?
	\item Quais limitações permanecem nas soluções existentes?
	\item Qual lacuna de conhecimento ou de tecnologia justifica a realização deste trabalho?
	\item Qual é exatamente a pergunta de pesquisa que o trabalho pretende responder?
	\item Qual é a hipótese, premissa ou proposição central que será investigada?
\end{enumerate}

{Resultado esperado.}
Ao final da contextualização, o leitor deve compreender
\textcolor{blue}{onde o problema está, por que ele importa e qual lacuna justifica o
	trabalho}.

% ---------------------------------------------------------
\section{Objetivos}
% ---------------------------------------------------------

\subsection{Objetivo geral}

\begin{enumerate}
	\item Qual é a finalidade científica ou tecnológica central do trabalho?
	\item Qual problema o trabalho pretende solucionar, analisar, modelar, desenvolver ou avaliar?
	\item Qual resultado concreto deverá existir ao final da pesquisa?
	\item O objetivo geral é compatível com o tempo, os recursos e o nível de um TCC?
	\item O objetivo é verificável por meio de experimentos, simulações, análise matemática ou métricas?
	\item O objetivo está diretamente relacionado à lacuna apresentada na contextualização?
\end{enumerate}

\subsubsection{Objetivos específicos}

\begin{enumerate}
	\setcounter{enumi}{6}
	\item Quais etapas são necessárias para alcançar o objetivo geral?
	\item É necessário realizar revisão bibliográfica específica?
	\item É necessário desenvolver um modelo matemático ou computacional?
	\item É necessário implementar, simular, testar, comparar ou validar uma solução?
\end{enumerate}

%Estrutura lógica recomendada.
%
%\begin{center}
%	\textbf{Problema $\rightarrow$ objetivo geral $\rightarrow$ objetivos específicos
%		$\rightarrow$ método $\rightarrow$ resultados $\rightarrow$ conclusão.}
%\end{center}
\textcolor{blue}{Note que isto será utilizado na seção de Análise e Discussões de Resultados e na seção de Conclusão
sobre seguinte princípio: O objetivo geral e os objetivos específicos aqui declarados foram cumpridos?
}

% ---------------------------------------------------------
\section{Contribuições}
% ---------------------------------------------------------

\begin{enumerate}
	\item Qual é a principal contribuição científica ou tecnológica do trabalho?
	\item O que foi efetivamente desenvolvido pelo autor?
	\item O que diferencia a solução proposta das existentes?
	\item O trabalho produz um novo modelo, método, algoritmo, arquitetura, dispositivo, sistema ou procedimento?
	\item Há contribuição teórica ou matemática?
	\item Há contribuição experimental ou computacional?
	\item Há contribuição metodológica?
	\item Há contribuição para uma aplicação prática ou industrial?
	\item Quais resultados podem ser reutilizados por outros pesquisadores?
	\item Qual conhecimento novo permanece após a conclusão do trabalho?
\end{enumerate}

\textcolor{blue}{
Pergunta-chave.
Se o trabalho fosse removido da literatura, \textbf{o que deixaria de existir?}
}

% ---------------------------------------------------------
\section{Publicações}
% ---------------------------------------------------------

\begin{enumerate}
	\item O trabalho resultou em artigo, congresso, periódico, capítulo ou outro produto científico?
	\item Qual parte do TCC originou cada publicação?
	\item Qual é a relação entre a publicação e os objetivos do TCC?
	\item Quais resultados foram apresentados na publicação?
	\item Houve diferenças entre a versão publicada e a versão apresentada no TCC?
	\item Quem são os autores da publicação?
	\item Qual é o veículo de publicação?
	\item Qual é a situação da publicação: submetida, aceita, publicada ou em preparação?
	\item Quais contribuições do TCC foram validadas ou disseminadas pela publicação?
	\item Existem direitos autorais, propriedade intelectual ou restrições associadas à publicação?
\end{enumerate}

% ---------------------------------------------------------
\section{Organização do restante do trabalho}
% ---------------------------------------------------------

\begin{enumerate}
	\item O que será apresentado no Capítulo~2?
	\item Por que o Estado da Arte é necessário para compreender o problema?
	\item O que será apresentado no Capítulo~3?
	\item Quais conceitos fundamentais precisam ser dominados pelo leitor?
	\item O que será apresentado no Capítulo~4?
	\item Como a solução proposta será formalizada?
	\item O que será apresentado no Capítulo~5?
	\item Como os resultados serão analisados e validados?
	\item Onde serão apresentadas as conclusões?
	\item Como os capítulos se conectam logicamente?
\end{enumerate}
% =========================================================
\chapter{Estado da Arte}
% =========================================================

O Estado da Arte (EA) deve identificar o conhecimento acumulado, as principais abordagens existentes, suas limitações e, sobretudo, a lacuna científica ou tecnológica na qual o trabalho está inserido.
Noutras palavras\textcolor{blue}{, o EA posiciona a solução de um TCC no espaço-tempo das soluções atuais, visando apresentar o que há de mais moderno e comparando-o com a solução ora proposta. }

Para fins de orientaçoes gerais na produção do TCC, deve-se trabalhar conforme as orientações disponibilizadas no sítio do IFCE:
\begin{enumerate}
    \item \textbf{Orientações Gerais para Entrega de Trabalho Final} \\ \url{https://acervo.ifce.edu.br/proen/bibliotecas/entrega-de-trabalhos-academicos}
    
    \item \textbf{PPC do curso de Engenharia de Controle e Automação}\\ \url{https://ifce.edu.br/maracanau/menu/cursos/superiores/bacharelados/controle_automacao/pdf/ppc-eca-1.pdf}
    
\end{enumerate}


% ---------------------------------------------------------
\section{Considerações iniciais}
% ---------------------------------------------------------

\begin{enumerate}
	\item Qual é o domínio científico específico investigado?
	\item Quais são os principais subtemas envolvidos?
	\item Quais critérios foram utilizados para selecionar os trabalhos analisados?
	\item Qual período da literatura foi considerado?
	\item Quais bases de dados foram utilizadas?
	\item Quais palavras-chave foram utilizadas nas buscas?
	\item Quais critérios de inclusão e exclusão foram adotados?
	\item A revisão é narrativa, sistemática, bibliométrica ou uma combinação dessas abordagens?
	\item Quais são as principais \textbf{tendências} identificadas?
	\item Como a revisão conduzirá à identificação da lacuna científica?
\end{enumerate}

% ---------------------------------------------------------
\section{Fronteira do conhecimento nos assuntos envolvidos}
% ---------------------------------------------------------

\begin{enumerate}
	\item Quais são os trabalhos mais relevantes e recentes sobre o problema?
	\item Quais pesquisadores ou grupos são referências no tema?
	\item Quais métodos são predominantes na literatura?
	\item Quais tecnologias ou arquiteturas são atualmente utilizadas?
	\item Quais modelos matemáticos aparecem com maior frequência?
	\item Quais métricas são utilizadas para avaliar as soluções?
	\item Quais resultados são considerados estado da arte?
	\item Quais limitações os próprios autores apontam?
	\item Quais resultados da literatura são conflitantes ou inconclusivos?
	\item Qual lacuna permanece aberta após a análise crítica da literatura?
\end{enumerate}

{Pergunta central.}
\textbf{Qual é o estado atual do conhecimento e exatamente onde meu trabalho se
	posiciona?}

% ---------------------------------------------------------
\section{Visão holística do estado da arte}
% ---------------------------------------------------------

\begin{enumerate}
	\item Como os diferentes trabalhos da literatura se relacionam?
	\item Quais são as principais linhas de pesquisa?
	\item Quais abordagens podem ser agrupadas por método?
	\item Quais abordagens podem ser agrupadas por aplicação?
	\item Quais variáveis ou métricas permitem comparar os trabalhos?
	\item Quais são as vantagens e desvantagens de cada abordagem?
	\item Quais soluções são complementares?
	\item Quais soluções são concorrentes?
	\item Como os trabalhos evoluíram historicamente?
	\item Qual espaço científico permanece disponível para a solução proposta?
\end{enumerate}

Sugestão de síntese comparativa.

\begin{center}
	\begin{tabular}{p{2.5cm}p{2cm}p{2cm}p{2cm}p{2cm}p{3cm}}
		\toprule
		\textbf{Trabalho} &
		\textbf{Método} &
		\textbf{Aplicação} &
		\textbf{Dados} &
		\textbf{Métrica} &
		\textbf{Principal resultado/limitação} \\
		\midrule
		Autor 1 & Método & Aplicação & Dados & Métrica & Resultado e limitação \\
		Autor 2 & Método & Aplicação & Dados & Métrica & Resultado e limitação \\
		Autor 3 & Método & Aplicação & Dados & Métrica & Resultado e limitação \\
		\bottomrule
	\end{tabular}
\end{center}

\textcolor{blue}{Note que isto traz destaque amplo para o que será proposto.}
% ---------------------------------------------------------
\section{Considerações finais}
% ---------------------------------------------------------

\begin{enumerate}
	\item Quais foram as principais conclusões do Estado da Arte?
	\item Quais abordagens apresentaram melhor desempenho?
	\item Quais limitações foram identificadas?
	\item Quais lacunas permanecem?
	\item Qual lacuna será atacada pelo presente trabalho?
	\item Por que essa lacuna é relevante?
	\item Como a solução proposta se posiciona em relação à literatura?
	\item Qual hipótese emerge da revisão?
	\item Quais requisitos a solução proposta deve atender?
	\item Como o Estado da Arte justifica o capítulo seguinte?
\end{enumerate}

\textcolor{blue}{O redator do TCC deve atentar a essas perguntas, pois a banca de avaliação haverá de realizá-las sobre o trabalho por ocasião da defesa.}

% =========================================================
\chapter{Fundamentação Teórica}
% =========================================================

A Fundamentação Teórica deve apresentar os conhecimentos indispensáveis para
que o leitor compreenda o problema, a solução proposta e os resultados
posteriormente apresentados.

% ---------------------------------------------------------
\section{Considerações iniciais}
% ---------------------------------------------------------

\begin{enumerate}
	\item Quais conceitos são indispensáveis para compreender o trabalho?
	\item Quais conhecimentos são pressupostos?
	\item Quais conceitos precisam ser formalmente definidos?
	\item Quais teorias sustentam a solução proposta?
	\item Quais variáveis e grandezas serão utilizadas posteriormente?
	\item Quais conceitos serão utilizados no modelo matemático?
	\item Quais conceitos serão utilizados nas simulações?
	\item Qual é a relação entre os fundamentos teóricos apresentados?
	\item O capítulo está fornecendo apenas o conhecimento necessário ou está se transformando em um tratado sobre o assunto?
	\item Como a fundamentação prepara o leitor para compreender a solução?
\end{enumerate}

% ---------------------------------------------------------
\section{Conteúdo seminal do assunto correlato}
% ---------------------------------------------------------

\begin{enumerate}
	\item Qual é a origem histórica do conceito ou tecnologia?
	\item Quem foram os pesquisadores responsáveis pelas contribuições fundamentais?
	\item Quais trabalhos são considerados seminais?
	\item Qual problema originalmente motivou o desenvolvimento da teoria?
	\item Quais são os princípios físicos, matemáticos ou computacionais envolvidos?
	\item Quais equações fundamentais descrevem o fenômeno?
	\item Quais hipóteses e simplificações estão presentes nos modelos?
	\item Quais são as limitações da teoria?
	\item Como a teoria evoluiu até o estado atual?
	\item Quais fundamentos serão efetivamente utilizados na solução proposta?
\end{enumerate}

{Observação.}
O termo ``seminal'' não significa necessariamente ``antigo''. O critério deve
ser a relevância fundamental do conteúdo para a compreensão do problema, e
não simplesmente sua antiguidade.

% ---------------------------------------------------------
\section{Considerações finais}
% ---------------------------------------------------------

\begin{enumerate}
	\item Quais fundamentos teóricos são essenciais para o trabalho?
	\item Quais equações serão utilizadas posteriormente?
	\item Quais conceitos serão utilizados no modelo?
	\item Quais parâmetros deverão ser definidos experimentalmente ou por simulação?
	\item Quais hipóteses teóricas serão adotadas?
	\item Quais limitações teóricas podem afetar os resultados?
	\item Como a fundamentação sustenta a solução proposta?
	\item O leitor possui, ao final deste capítulo, conhecimento suficiente para compreender o próximo?
\end{enumerate}
% =========================================================
\chapter{Solução Proposta}
% =========================================================

Este capítulo deve responder essencialmente à pergunta:

\begin{center}
	\textbf{O que foi desenvolvido e como funciona?}
\end{center}

% ---------------------------------------------------------
\section{Arquitetura holística do sistema}
% ---------------------------------------------------------

\begin{enumerate}
	\item Qual é a arquitetura geral da solução?
	\item Quais são os principais componentes?
	\item Qual é a função de cada componente?
	\item Como os componentes se comunicam?
	\item Quais são as entradas do sistema?
	\item Quais são os processos internos?
	\item Quais são as saídas?
	\item Quais interfaces existem entre os subsistemas?
	\item Quais componentes são físicos, computacionais, matemáticos ou experimentais?
	\item Como a arquitetura proposta se diferencia das arquiteturas existentes?
\end{enumerate}

{Recomendação.}
Apresentar um diagrama de blocos geral antes de entrar nos detalhes dos
componentes individuais.

% ---------------------------------------------------------
\section{Modelo matemático do sistema}
% ---------------------------------------------------------

\begin{enumerate}
	\item Qual fenômeno ou comportamento está sendo modelado?
	\item Quais são as variáveis de entrada?
	\item Quais são as variáveis de saída?
	\item Quais são os parâmetros do sistema?
	\item Quais leis físicas ou princípios matemáticos governam o sistema?
	\item Quais equações descrevem o comportamento do sistema?
	\item Quais hipóteses foram adotadas?
	\item Quais simplificações foram necessárias?
	\item Quais condições iniciais e de contorno foram consideradas?
	\item O modelo pode ser validado experimentalmente ou comparado com modelos conhecidos?
\end{enumerate}

{Sequência recomendada.}

\begin{center}
	\textbf{
		Fenômeno físico
		$\rightarrow$
		hipóteses
		$\rightarrow$
		variáveis
		$\rightarrow$
		equações
		$\rightarrow$
		parâmetros
		$\rightarrow$
		condições
		$\rightarrow$
		modelo final.
	}
\end{center}

% ---------------------------------------------------------
\section{Fluxograma do sistema}
% ---------------------------------------------------------

\begin{enumerate}
	\item Qual é o estado inicial do sistema?
	\item Qual é o primeiro processamento realizado?
	\item Quais decisões precisam ser tomadas?
	\item Quais condições determinam cada caminho?
	\item Existem ciclos ou realimentações?
	\item Existem processos iterativos?
	\item Quais dados entram em cada etapa?
	\item Quais dados saem de cada etapa?
	\item Qual é o critério de parada?
	\item Qual é o resultado final do algoritmo ou processo?
\end{enumerate}

O fluxograma deve permitir que um terceiro pesquisador reconstrua a lógica
operacional da solução.

% ---------------------------------------------------------
\section{Considerações finais}
% ---------------------------------------------------------

\begin{enumerate}
	\item Qual solução foi efetivamente desenvolvida?
	\item Quais são seus principais componentes?
	\item Quais são suas principais características?
	\item Quais hipóteses foram adotadas?
	\item Quais são suas limitações?
	\item Quais parâmetros serão avaliados?
	\item Como a solução será testada?
	\item Quais critérios serão utilizados para determinar se ela funciona?
	\item Como ela será comparada às soluções da literatura?
	\item O que precisa ser demonstrado experimentalmente ou por simulação?
\end{enumerate}


% =========================================================
\chapter{Simulações, Análises e Discussão de Resultados}
% =========================================================

Este capítulo deve apresentar evidências capazes de sustentar as afirmações
científicas do trabalho. Resultado não deve ser tratado simplesmente como uma
figura, tabela ou valor numérico: cada resultado precisa ser interpretado.

% ---------------------------------------------------------
\section{Considerações iniciais}
% ---------------------------------------------------------

\begin{enumerate}
	\item Qual é a finalidade dos experimentos ou simulações?
	\item Quais objetivos específicos estão sendo avaliados?
	\item Quais hipóteses estão sendo testadas?
	\item Quais métricas serão utilizadas?
	\item Quais resultados são esperados?
	\item Como os resultados serão comparados?
	\item Qual é o critério de sucesso?
	\item Como a validade dos resultados será assegurada?
	\item Quais fontes de incerteza existem?
	\item Como os resultados serão discutidos?
\end{enumerate}

% ---------------------------------------------------------
\section{Cenários, métodos e parâmetros de simulação/teste/coleta de dados}
% ---------------------------------------------------------

\begin{enumerate}
	\item Qual ambiente foi utilizado?
	\item Qual software, hardware ou plataforma foi utilizada?
	\item Qual versão das ferramentas foi empregada?
	\item Quais parâmetros foram utilizados?
	\item Quais valores foram mantidos constantes?
	\item Quais parâmetros foram variados?
	\item Quais cenários foram considerados?
	\item Quantas repetições foram realizadas?
	\item Como os dados foram coletados e tratados?
	\item Como outro pesquisador poderia reproduzir o experimento?
\end{enumerate}

{Pergunta fundamental.}
\textbf{Um pesquisador externo conseguiria reproduzir exatamente este
	experimento apenas lendo esta seção?}

Se a resposta for negativa, a metodologia ainda não está suficientemente
documentada.

% ---------------------------------------------------------
\section{Apresentação e discussão dos resultados}
% ---------------------------------------------------------

Para cada resultado relevante, o estudante deve responder às seguintes
perguntas:

\begin{enumerate}
	\item O que está sendo mostrado?
	\item Por que esse resultado é importante?
	\item Qual variável está sendo analisada?
	\item Qual tendência pode ser observada?
	\item Por que essa tendência ocorre?
	\item O resultado confirma ou contradiz a hipótese?
	\item O resultado está de acordo com a literatura?
	\item Por que o resultado difere ou coincide com trabalhos anteriores?
	\item Qual é a magnitude da melhoria ou degradação observada?
	\item Qual conclusão científica pode ser extraída desse resultado?
\end{enumerate}

{Estrutura recomendada para discussão.}

\begin{center}
	\textbf{
		Resultado
		$\rightarrow$
		comportamento observado
		$\rightarrow$
		explicação física/matemática
		$\rightarrow$
		comparação
		$\rightarrow$
		interpretação
		$\rightarrow$
		conclusão.
	}
\end{center}

% ---------------------------------------------------------
\section{Considerações finais}
% ---------------------------------------------------------

\begin{enumerate}
	\item Quais foram os principais resultados?
	\item Quais objetivos específicos foram atendidos?
	\item Quais hipóteses foram confirmadas?
	\item Quais hipóteses foram rejeitadas?
	\item Qual foi o melhor cenário?
	\item Quais parâmetros mais influenciaram o desempenho?
	\item Quais limitações foram identificadas?
	\item Os resultados são estatisticamente ou tecnicamente significativos?
	\item Os resultados são coerentes com a literatura?
	\item Quais evidências permitem afirmar que a solução funciona?
\end{enumerate}


% =========================================================
\chapter{Conclusões e Trabalhos Futuros}
% =========================================================

\section{Conclusões}

\begin{enumerate}
	\item Qual era o problema investigado?
	\item Qual era o objetivo geral?
	\item O objetivo geral foi alcançado?
	\item Quais objetivos específicos foram alcançados?
	\item Quais foram as principais descobertas?
	\item Qual é a principal contribuição científica ou tecnológica?
	\item Quais resultados mais relevantes sustentam a conclusão?
	\item Quais limitações permaneceram?
	\item Quais questões permaneceram sem resposta?
	\item Qual é a principal mensagem científica que deve permanecer após a leitura do trabalho?
\end{enumerate}

\section{Trabalhos futuros}

\begin{enumerate}
	\item Que limitações poderiam ser eliminadas?
	\item Que novos cenários poderiam ser investigados?
	\item O método pode ser generalizado?
	\item O sistema pode ser implementado em escala real?
	\item O método pode ser comparado com novas tecnologias?
	\item Quais novos modelos poderiam ser incorporados?
	\item Quais testes experimentais ainda seriam necessários?
	\item Existe potencial para publicação científica ou patente?
	\item Quais aplicações práticas poderiam ser exploradas?
	\item Qual seria a evolução natural deste TCC para uma dissertação de mestrado ou projeto de pesquisa?
\end{enumerate}






% =========================================================
\section{Referências}
% =========================================================

Embora as referências não constituam propriamente um capítulo argumentativo,
a construção da bibliografia também pode ser orientada por perguntas.

Note que, utilizando o \LaTeX, todas as referências deve estar contidas no arquivo \texttt{Referencias.bib}.Sendo, evocadas no TCC pelos comandos abaixo, dentre outros: 
\texttt{/\/cite{}, /\/citeonline{},/\/citeauthor{bibid},/\/citeauthoronline{}}.


\begin{enumerate}
	\item Todas as afirmações científicas relevantes estão devidamente referenciadas?
	\item As referências utilizadas são primárias sempre que possível?
	\item Os trabalhos mais relevantes e recentes foram considerados?
	\item Os trabalhos seminais foram identificados?
	\item As referências são provenientes de fontes academicamente confiáveis?
	\item Os artigos citados realmente sustentam as afirmações feitas?
	\item Há excesso de referências secundárias?
	\item Há referências antigas que deveriam ser complementadas por trabalhos recentes?
	\item Todas as referências citadas no texto aparecem na lista final?
	\item Todas as referências listadas são efetivamente citadas no texto?
\end{enumerate}

Veja o exemplo a seguir:

``
...
A análise da estabilidade e do comportamento transitório em malha fechada é um passo crítico no desenvolvimento de sistemas de automação industrial. Para mitigar erros em regime permanente sem comprometer a velocidade de resposta do sistema, a literatura técnica frequentemente recorre à sintonia de controladores Proporcional-Integral-Derivativo (PID). 

Conforme aponta \citeonline{ogata1998engenharia}, a ação integradora elimina o erro residual frente a pertubações em degrau, enquanto a ação derivativa antecipa a tendência do erro, adicionando amortecimento ao sistema. Desse modo, o projeto de compensadores baseados no método do lugar das raízes \cite{ogata1998engenharia} permite posicionar os polos dominantes de forma a satisfazer os requisitos de sobressinal e tempo de acomodação predefinidos.

...

''
% =========================================================
\chapter{Apêndices}
% =========================================================

O apêndice deve conter material produzido pelo próprio autor que seja
relevante para complementar ou permitir a reprodução do trabalho, mas que,
por sua extensão ou natureza, não seja adequado ao corpo principal do texto.

\begin{enumerate}
	\item Existe material produzido pelo autor que complementa o trabalho?
	\item Esse material é excessivamente extenso para permanecer no corpo do texto?
	\item Ele é importante para a reprodução do trabalho?
	\item Há códigos-fonte relevantes?
	\item Há tabelas extensas?
	\item Há deduções matemáticas detalhadas?
	\item Há questionários ou instrumentos de coleta desenvolvidos pelo autor?
	\item Há parâmetros completos de simulação?
	\item Há procedimentos experimentais detalhados?
	\item O leitor pode consultar esse material para aprofundar ou reproduzir o trabalho?
\end{enumerate}

{Exemplos de conteúdo de apêndices:}

\begin{itemize}
	\item Código desenvolvido;
	\item Dedução matemática extensa;
	\item Tabelas completas;
	\item Parâmetros de simulação;
	\item Questionários;
	\item \textit{Scripts};
	\item Procedimentos experimentais.
\end{itemize}
% =========================================================
\chapter{Anexos}
% =========================================================

O anexo deve conter material produzido por terceiros ou documentação externa
que seja necessária ou relevante para complementar o trabalho.

\begin{enumerate}
	\item Existe documentação externa necessária para compreender o trabalho?
	\item O material foi produzido por outro autor ou instituição?
	\item Há normas técnicas relevantes?
	\item Há \textit{datasheets}?
	\item Há manuais?
	\item Há especificações de fabricantes?
	\item Há documentos oficiais?
	\item Há legislação ou regulamentação pertinente?
	\item Há instrumentos ou documentos externos utilizados na pesquisa?
	\item O documento realmente precisa estar anexado ou uma referência seria suficiente?
\end{enumerate}



% Final de numeração das páginas.
\pagestyle{empty}

\bibliography{3-PosTextuais/Referencias.bib} \thispagestyle{empty}

% Indicação de inicialização do apendice.
\appendix

%
%%%%%%%%%%%%%%%%%%%%%%%%%%%%%%%%%%%%%%%%%%%%%%%%%%%%%%%%%%%%%%%%%%%%%
%%%%%%%%%%%%%%%%%%%%%%%%%%%% Apendice %%%%%%%%%%%%%%%%%%%%%%%%%%%%%%%
%%%%%%%%%%%%%%%%%%%%%%%%%%%%%%%%%%%%%%%%%%%%%%%%%%%%%%%%%%%%%%%%%%%%%


% Inicio do Apendice
\begin{apendicesenv}


\chapter{Referências nos Arquivos .bib}\label{ApeExe}
\thispagestyle{empty}


\begin{description}
	\item[Para artigos em revistas ]:
		\small
		\begin{verbatim}
			@ARTICLE{artigo,
  		  AUTHOR = {Nome do autor and Nome do Outro Autor},
  		  TITLE = {Título do Artigo},
  		  JOURNAL = {Nome da revista em que o artigo foi publicado},
  		  YEAR = {Ano},
  		  volume = {Volume},
  		  number = {Número},
  		  pages = {pp.inicial-pp.final},
  		  month = {Mês},
  		  address = {Local},
			}
		\end{verbatim}

	\item[Livros]:
		\small
		\begin{verbatim}
			@BOOK{livros,
  			AUTHOR = {Nome Completo do autor1 and Nome Completo do autor2},
  			editor = {Quem organizou o livro},
  			TITLE = {Título do livro},
  			PUBLISHER = {Editora},
  			YEAR = {Ano},
  			volume = {Volume},
  			number = {Número},
  			series = {Série a qual o livro pertence},
  			address = {Local},
  			edition = {Número da edição},
  			month = {mês},
			}
		\end{verbatim}

	\item[Capítulos em Livros ]:
		\small
		\begin{verbatim}
			@INBOOK{ComoVoceIraReferenciarNoTextoEstaReferencia,
  			AUTHOR = {Nome Completo do autor1 and Nome Completo do autor2},
  			editor = {Quem organizou o livro},
  			TITLE = {Título do Livro},
  			CHAPTER = {Título do Capítulo},
  			pages = {pp.inicial-pp.final},
  			PUBLISHER = {Editora},
  			YEAR = {Ano},
  			volume = {Volume},	
  			number = {Número}
			}
		\end{verbatim}
		
	
\end{description}


\end{apendicesenv}
%%%%%%%%%%%%%%%%%%%%%%%%%%%%%%%%%%%%%%%%%%%%%%%%%%%%%%%%%%%%%%%%%%%%%
%%%%%%%%%%%%%%%%%%%%%%%%%%%% Anexo %%%%%%%%%%%%%%%%%%%%%%%%%%%%%%%%%%
%%%%%%%%%%%%%%%%%%%%%%%%%%%%%%%%%%%%%%%%%%%%%%%%%%%%%%%%%%%%%%%%%%%%%


% Inicio do Anexo
\begin{anexosenv}

\chapter{História da Linguagem Tex}\label{AxeExe}
\thispagestyle{empty}

Quando o primeiro volume de The Art of Computer Programming de Donald Knuth foi publicado em 1968, o livro tinha sido tipografado com uma composição de metal quente da Monotype Corporation. Este método datado do século XIX tinha um bom estilo clássico, que agrava Knuth. Quando a segunda edição foi publicada em 1976, todo o trabalho precisou ser tipografado novamente porque a tecnologia da Monotype Corporation tinha sido substituída pela fotocomposição e as fontes originais não estavam mais disponíveis. Quando Knuth recebeu as provas da nova versão em 30 de março de 1977, ele as achou terríveis. Naquela época, Knuth viu pela primeira vez a saída de um sistema de tipografia digital de alta qualidade e se interessou pela tipografia digital. O desapontamento com as provas deram a Knuth a motivação final para resolver o problema de uma vez por todas, com o desenvolvimento do seu próprio sistema de tipografia. Em 13 de maio de 1977, ele escreveu um memorando para si mesmo descrevendo as características básicas do TeX.



\end{anexosenv}

\end{document}
